% Latex Document Class used exclusively by Engenius UA 
% (c) 2018, Diogo Correia and João Santos, MIT Licensed
% Repository: https://github.com/dvcorreia/engenius-ua-latex-template.git
%

% Put between [ ] in document class your department following these macros:
%
% + ec - Electronics and Communications Department
% + epower - Powertrain Electric Department
% + business - Business Department
% + cpower - Powertrain Combustion Department
% + svd - Suspension and vehicle dynamics Department
% + chassis - Chassis develpment Department
%

\documentclass[]{engenius}

% Fill the document parameters
% Remove \titletwo if you want one title
% You can remove authors by deleting then
% You have 3 option for the department identification as showed bellow

\title{Engenius Latex Template}
\titletwo{How to use it!}

\author{Author 1 \\ \email{emailauthor1@ua.pt}
	\and Author 2 \\ \email{emailauthor2@ua.pt} \\ \depofepower
    \and Author 3 \\ \email{emailauthor3@ua.pt} \\ \depofbusinesscolor
    } 
    
% Update date for the release date

% For versions: #1.#2.#3 where:
% #1 represents a complete modification of the previous work;
% #2 a big change in the current work;
% #3 small changes and updates done to the current work.

\date{\today}
\version{0.1.0}

\begin{document}

\maketitle

\begin{abstract}

The \textbf{abstract} is supposed to be a brief resume where you write about what the report will be about, the objectives and if is an update to a older report from you or other author (that must be referenced here and in the Revisions section). An introduction is optional since the main objective is a compact study and discussion of the development process.

\end{abstract}

\section{How to write well}
In this template we created some commands that provide to a non \LaTeX~ user an easy  write without any knowledge about \TeX~or~\LaTeX.

\subsection{Document Information}
Before start writing you should complete the preamble with your information. Don't forget to insert your department identification in \mintinline{latex}{\documentclass[your department]{engenius}} as told in the comments.

Fill the authors information and choose a common way of identifying your department. There are 3 options you can choose: not identifying at all, identifying as showed in Author2 (\mintinline{latex}{\depofepower})or with your department color and exemplified in Author3 (\mintinline{latex}{\depofbusinesscolor}).

Versions are important since they can evaluate the progress of the work. A version is composed by 3 integers, \textit{version N1.N2.N3}. N1 increases when a complete modification of the previous work is implemented, N2 increases when are made big changes in the current work, N3 increases when small changes like updates are done or added.

\subsection{Writing the document}

Next will be indicated good practices when writting with this template.

\subsubsection{Equations}

To use equations you can make them inline like this $x = y + z$ or you can do them like this:

\begin{equation}\label{eq:timehtol}
    \int_{0}^{tp_{HL}}dt = tp_{HL} = -C_L \int_{V_{dd}}^{V_{dd}/2}\frac{1}{I_{DSN}} dV_O
\end{equation}

And you can reference it like this \mintinline{Latex}{\ref{eq:timehtol}} , equation~\ref{eq:timehtol}. It will automatically his change if you put other equations before, so this way you dont need to worry about equation numbering.

\subsubsection{Images}

Images are placed like the example bellow. You can tune the width to make it fit your needs, it can take any units (cm, in, em ...). To reference an image you can follow the same aproach as in equations: \mintinline{Latex}{\ref{img:engeniuslogo}}, image~\ref{img:engeniuslogo}.

\begin{figure}[!ht]
    \begin{center}
        \includegraphics[width=\textwidth/3]{src/engeniusLogo.png}
        \caption{Engenius Image Example}
        \label{img:engeniuslogo}
    \end{center}
\end{figure}

\subsubsection{Code}

Raw code can be inserted inline like this \mintinline{python}{print("Hello world!")}, or in a box like the one bellow. Check our code to see how to do it.

\begin{codebox}{Example Code}
    \begin{minted}{C}
        #include <stdio.h>
        int main()
        {
           // printf() displays the string inside quotation
           printf("Hello, World!");
           return 0;
        }
    \end{minted}
\end{codebox}

\subsubsection{References and Revision}

References must be used when you use information from external fonts. There is no shame in referencing everything :') . Your bibliografy must be in the file \mintinline{text}{./biblio.bib} structured like the examples already there. To cyte something just use \mintinline{Latex}{\cite{name}}, like this \cite{einstein}. The reference will then show in the references section.

The revision section is used to tell which report are you updating/revising in case you have one. This section exist in order to give credits to the people that worked before you, they must not be forgotten :) . You also need to put in your document all revisions that were in the document you are updating and update the version according.

\section{Questions, Bugs and Suggestions}

This template was created by Diogo Correia and João Santos.
Original Repository \textcolor{blue}{\href{https://github.com/dvcorreia/engenius-ua-latex-template}{here}}.

We spent a good amount of time creating this template and we are open for collaborators and suggestions :) .

If you find any bugs please make us know so we can fix then.
If you have any question or need help you can also contact us.

All of that can be done in github in the following link:  \textcolor{blue}{\href{https://github.com/dvcorreia/engenius-ua-latex-template/issues/new}{here}}, by selecting the correct label (bugs for bugs, enhancement for suggestions and  questions with the question label).

Thank for the attention and continuation of a good work :) .



\bibliography{biblio}

\revisions
\revisionof{Engenius Latex Template: How to use it!}{Electronics/Communications}{Diogo Correia, João Santos}{0.0.1}{12/02/2018}

\end{document}
